\newcommand{\lrparen}[1]{% \lrparen{..}
  \left(#1\right)}
\newcommand{\lrcurly}[1]{
  \left\{#1\right\}}
\newcommand{\lrbracket}[1]{
  \left[#1\right]}
\newcommand{\lrceil}[1]{
  \left\lceil#1\right\rceil}
\newcommand{\llrrbracket}[1]{
  \left[\mkern-3mu\left[#1\right]\mkern-3mu\right]}
\newcommand{\llrrparen}[1]{
  \left(\mkern-6mu\left(#1\right)\mkern-6mu\right)}
\newcommand{\lrangle}[1]{
  \left\langle#1\right\rangle}
\newcommand{\dan}[1]{
\textcolor{red}{(Dan) #1}}
\newcommand{\remy}[1]{
\textcolor{blue}{(Remy) #1}}
\definecolor{Gray}{gray}{0.8}

\newcommand{\Bold}[0]{B_{old}}
\newcommand{\Bcur}[0]{B_{cur}}
\newcommand{\Bnew}[0]{B_{new}}
\newcommand{\MPLAN}[0]{\textsc{LA}}
\newcommand{\RPLAN}[0]{\textsc{RA}}
\newcommand{\rmul}[0]{*}
\newcommand{\radd}[0]{+}

\newcommand{\unbind}[2]{{\ensuremath{\scriptstyle\left[-#1,-#2\right]\,}}}
\newcommand{\bind}[1]{{\ensuremath{\scriptstyle\left[#1\right]\,}}}
\newcommand{\tr}[0]{{\mathsf{T}}}
\newcommand{\canon}{\mathcal C}

\newtheorem{definition}{Definition}[section]
\newtheorem{theorem}{Theorem}[section]
\newtheorem{lemma}[theorem]{Lemma}
\newtheorem{example}{Example}
\newtheorem{corollary}{Corollary}


\newcommand{\sys}{SPORES}


\newcommand{\Z}{\mathbb Z} % integers
\newcommand{\N}{\mathbb N} % the natural numbers
\newcommand{\R}{\mathbb R} % the real numbers
\newcommand{\Q}{\mathbb Q} % the rational numbers
\newcommand{\C}{\mathbb C} % complex numbers
\newcommand{\D}{\mathbf D} % bold-face D, used for generic domain

